% Generated draft: requires mathematical and source review.
\documentclass[11pt]{article}
\usepackage[margin=25mm]{geometry}
\usepackage{fontspec}
\setmainfont{Latin Modern Roman}
\usepackage{amsmath,amssymb,mathtools,unicode-math}
\setmathfont{Latin Modern Math}
\usepackage{enumitem,needspace,xurl,hyperref}
\setlength{\parindent}{0pt}
\setlength{\parskip}{5pt}
\setlength{\emergencystretch}{3em}
\newcommand{\N}{\mathbb{N}}
\newcommand{\Z}{\mathbb{Z}}
\newcommand{\Q}{\mathbb{Q}}
\newcommand{\R}{\mathbb{R}}
\newcommand{\C}{\mathbb{C}}
\newcommand{\F}{\mathbb{F}}
\newcommand{\A}{\mathbb{A}}
\newcommand{\PP}{\mathbb{P}}
\newcommand{\E}{\mathbb{E}}
\newcommand{\eps}{\varepsilon}
\newcommand{\dd}{\,\mathrm{d}}
\begin{document}
\section*{UnsolvedMath ID: 800524}
\textbf{Source problem code:} catalogue entry
\subsection{Definitions and mathematical statement}
Let \(A\in\R^{n\times d}\) be full rank, let \(B\in\R^{n\times p}\), and let \(\Omega\in\R^{n\times k}\). The source requires \(\Omega\) to be an ``oblivious subspace injection with dimension \(d\) and injectivity \(1-\varepsilon\).'' This condition is retained verbatim because the supplied record does not define it. In particular, the sketching map acting on the data is \(\Omega^\top:\R^n\to\R^k\).

For a real matrix \(M=(M_{ij})\), define its Frobenius norm by \(\|M\|_{\mathrm F}=(\sum_{i,j}|M_{ij}|^2)^{1/2}\), and write \(M^+\) for its Moore--Penrose pseudoinverse. Set
\[
X=A^+B\in\R^{d\times p},\qquad
\widetilde X=(\Omega^\top A)^+(\Omega^\top B)\in\R^{d\times p}.
\]
Thus \(X\) is an ordinary least-squares solution, and \(\widetilde X\) is the sketch-and-solve approximation specified by the source.

\textbf{Question (source Problem 33).} Prove or disprove that, for every choice of matrices satisfying the stated hypotheses, one necessarily has
\[
\|A\widetilde X-B\|_{\mathrm F}
\leq \bigl(1+\mathcal O(\varepsilon)\bigr)
\min_{Y\in\R^{d\times p}}\|AY-B\|_{\mathrm F}.
\]
The admissible range of \(\varepsilon\), the asymptotic regime and permitted dependencies of the implicit constant, and any probability qualification are unspecified in the supplied record. It also supplies no relation between \(n\) and \(d\). Consequently, this is a faithful statement of the source question with unresolved qualifications, rather than a fully specified quantified conjecture.
\subsection{Short English statement}
Does sketch-and-solve achieve a Frobenius residual within a factor \(1+\mathcal O(\varepsilon)\) of the optimal least-squares residual under the source's stated oblivious subspace injection hypothesis? The meaning of that hypothesis, the probabilistic guarantee, and the uniformity of the error factor require clarification.
\subsection{Sources}
\begin{itemize}
\item Local UnsolvedMath registry, permanent ID \(800524\), source code \texttt{catalogue entry}, title ``Linear Systems and Eigenvalue Problems --- Problem 33 (Sketch-and-solve).'' The supplied registry labels the problem open; this label is not an independently verified mathematical or literature assessment.
\item Amsel, Baumann et al., \emph{Linear Systems and Eigenvalue Problems: Open Questions from a Simons Workshop} (2026), Problem 33. Supplied source URL: \url{https://arxiv.org/abs/2602.05394}. The registry attributes the problem to Amsel, Baumann et al., with proposed year 2026 and attribution/context ``Sketch-and-solve.''
\item Supplied provenance records access on 2026-10-16 and extraction by \texttt{arxiv\_latex}. This access date is source metadata; no independent source visit or verification is claimed. The supplied upstream revision is \texttt{d371ee917ac9647fbb06cd6fc529deec6f3dda82}.
\end{itemize}
\end{document}
